% Pandoc template for the fed_fraud documents.
% Based on proposal_gen/templates/detailed_proposal_template/detailed_1_classic.tex
% (Classic report): centered title page, table of contents, ruled section
% headings and the Geodesic house footer. Fonts are Latin Modern (lmodern).
% Two layouts: the full report (title page and contents) and, with
% -V brief, a compact one-page-brief header in the same style.
\documentclass[letterpaper,11pt]{article}
\usepackage[T1]{fontenc}
\usepackage[utf8]{inputenc}
\usepackage{lmodern}
\usepackage[letterpaper,top=0.9in,bottom=1.2in,left=1in,right=1in,footskip=0.6in]{geometry}
\usepackage[table]{xcolor}
\usepackage{geodesicbrand}
\usepackage{amsmath,amssymb}
\usepackage{enumitem}
\usepackage{array}
\usepackage{calc}
\usepackage{longtable}
\usepackage{makecell}
\usepackage[most]{tcolorbox}
\usepackage{titlesec}
\usepackage{needspace}
\usepackage{etoolbox}
\usepackage{hyperref}
\hypersetup{colorlinks=true,urlcolor=coolblack,linkcolor=coolblack,filecolor=geoblue,
  pdftitle={Live demo},pdfauthor={\geofirm}}
\urlstyle{same}

\setlist{itemsep=2pt,leftmargin=1.4em}
\setlength{\parindent}{0pt}
\setlength{\parskip}{6pt}
\providecommand{\tightlist}{\setlength{\itemsep}{0pt}\setlength{\parskip}{0pt}}
\setcounter{secnumdepth}{0}

% Section style: house rule under each heading
\titleformat{\section}{\Large\bfseries\color{coolblack}}{\thesection}{0.6em}{}[{\color{coolblack}\titlerule[0.5pt]}]
\titleformat{\subsection}{\large\bfseries\color{geomid}}{\thesubsection}{0.6em}{}
\titleformat{\subsubsection}{\normalsize\bfseries\color{geomid}}{\thesubsubsection}{0.6em}{}
\titlespacing*{\section}{0pt}{10pt}{6pt}

% Tables are written by latex/filters.lua (navy header, banded rows)
\setlength{\LTpre}{6pt}\setlength{\LTpost}{6pt}

\begin{document}

% BRIEF HEADER
\thispagestyle{geodesic}
\begin{minipage}[c]{0.13\textwidth}\includegraphics[width=\linewidth]{icon.png}\end{minipage}\hfill
\begin{minipage}[c]{0.84\textwidth}
  {\textsc{\large \geofirm}}\par\vspace{2pt}
  {\small\color{geogray}Federal fraud expansion\enspace\textperiodcentered\enspace When: \compiletime}
\end{minipage}\par\vspace{6pt}
{\color{coolblack}\rule{\linewidth}{0.3mm}}\par\vspace{2pt}
{\LARGE\bfseries\color{coolblack} Live demo}\par
\vspace{2pt}{\large\color{geomid} The methods on real public data}\par
\vspace{-2pt}{\color{coolblack}\rule{\linewidth}{0.3mm}}\par\vspace{4pt}
\pagestyle{geodesic}

\textbf{Status:} Draft v0.1 (2026-11-04), for Steve's review. Extension
8 of 10 (see the \href{../extensions.pdf}{extensions overview}). Code:
\texttt{analytics/demo/run\_demo.py} (scores) and
\texttt{analytics/demo/fetch\_usaspending.py} (pulls awards); results:
\texttt{docs/demo/data/demo.json}; shown on the site at \texttt{/demo/}.

\hypertarget{bottom-line}{%
\section*{Bottom line}\label{bottom-line}}
\addcontentsline{toc}{section}{Bottom line}

The site's tool pages were wireframes with sample records. The demo
replaces ``imagine this'' with ``here it is'': the same reference code
that the forensic methods handbook describes, run on real public data,
with the results on a page. It deliberately starts with a dataset where
the honest answer is ``nothing to see'', because the most important
thing to show a leader is that the tests know their own limits.

What runs today, on data committed to the repo:

\begin{enumerate}
\def\labelenumi{\arabic{enumi}.}
\tightlist
\item
  \textbf{Digit test on every agency's FY2026 outlays, budget authority
  and obligations} (USAspending). 301 amounts from 101 agencies.
\item
  \textbf{Concentration} of federal outlays across agencies.
\item
  \textbf{Reconciliation:} agencies whose obligations or outlays exceed
  budget authority by more than 5\%.
\item
  \textbf{The ROI model} on GAO's FY2025 improper payment figures.
\end{enumerate}

What runs where the network allows (the ``Demo data'' GitHub Actions
workflow, or a laptop): contract awards for any agency from the
USAspending API, scored with the digit test, a round-number test,
bunching below the simplified acquisition threshold and recipient
concentration.

\hypertarget{results-on-agency-totals}{%
\section*{1. Results on agency totals}\label{results-on-agency-totals}}
\addcontentsline{toc}{section}{1. Results on agency totals}

\begingroup\small\setlength{\tabcolsep}{4pt}\renewcommand{\arraystretch}{1.25}
\rowcolors{2}{white}{geolight}
\begin{longtable}[l]{@{}>{\raggedright\arraybackslash}p{\dimexpr 0.302\linewidth-2\tabcolsep}>{\raggedright\arraybackslash}p{\dimexpr 0.349\linewidth-2\tabcolsep}>{\raggedright\arraybackslash}p{\dimexpr 0.349\linewidth-2\tabcolsep}@{}}
\rowcolor{coolblack}\leavevmode\color{white}\bfseries Test & \leavevmode\color{white}\bfseries Result & \leavevmode\color{white}\bfseries Reading \\
\endhead
First-digit test, all three amounts pooled (n = 301) & Chi-square p = 0.024; MAD 0.0195 (``nonconformity'' on Nigrini's scale); spike at digit 3 & \textbf{A false positive, and a useful one.} The three amounts per agency are near-copies of each other, so the 301 values are really about 100 independent numbers counted three times; that inflates the chi-square. Benford tests assume many independent transactions \\
First-digit test, outlays only (n = 101) & Chi-square p = 0.20; MAD 0.031 & No evidence against Benford. The MAD still reads ``nonconformity'' because Nigrini's MAD ranges assume thousands of records; at n = 101, MAD is mostly noise. This is why the handbook sets a minimum of 1,000 records \\
Concentration & Gini 0.94; top five agencies 85\% of outlays, top ten 94\% & Expected: Social Security, health, Treasury and defense dominate. Context for the scorecard, not a fraud signal \\
Reconciliation & No agency with more than \$100M in budget authority shows obligations or outlays more than 5\% above it & Clean at this level. Agency-level totals are reconciled before publication; the test earns its keep at the program and award level \\
\end{longtable}\endgroup

The page shows the observed and expected digit distributions side by
side so a viewer can see what a ``pass'' looks like before seeing a
``flag''.

\hypertarget{award-level-demo-when-run-with-network-access}{%
\section*{2. Award-level demo (when run with network
access)}\label{award-level-demo-when-run-with-network-access}}
\addcontentsline{toc}{section}{2. Award-level demo (when run with
network access)}

\texttt{fetch\_usaspending.py} pulls the most recent contract awards for
one agency and year (sorted by date, not size, so the sample is not only
the largest awards). \texttt{run\_demo.py} then reports, for each file:

\begin{itemize}
\tightlist
\item
  first-digit test (when there are at least 100 awards);
\item
  share of awards that are round thousands;
\item
  bunching below the simplified acquisition threshold (when at least 200
  awards fall between half and one and a half times it);
\item
  number of recipients, Herfindahl index and the top-ten recipient
  share.
\end{itemize}

No recipient is named on the page. Award-level results are statistics on
public records, for showing the method, not leads.

\hypertarget{how-to-refresh}{%
\section*{3. How to refresh}\label{how-to-refresh}}
\addcontentsline{toc}{section}{3. How to refresh}

\begin{verbatim}
python3 analytics/roi/model.py
python3 analytics/demo/run_demo.py
# with network access to api.usaspending.gov:
python3 analytics/demo/fetch_usaspending.py \
  --agency "Department of Veterans Affairs" --fy 2025
python3 analytics/demo/run_demo.py
\end{verbatim}

Or run the \textbf{Demo data} workflow from the repository's Actions
tab; it uploads \texttt{docs/demo/data/} as an artifact to commit.

\hypertarget{next-demos}{%
\section*{4. Next demos}\label{next-demos}}
\addcontentsline{toc}{section}{4. Next demos}

\begin{enumerate}
\def\labelenumi{\arabic{enumi}.}
\tightlist
\item
  CMS hospice enrollment and owner files through the network scorer
  (HHS-1), once the public files are pulled.
\item
  SAM.gov exclusions matched to USAspending recipients (cross-agency
  matching, business side).
\item
  Single Audit repeat findings for the largest pass-through programs.
\end{enumerate}

\end{document}
